\chapter{optional、variant、expected 风格和标准库词汇类型}

\section{问题入口：返回值为什么不能只靠特殊值}

查找用户时，找不到用户不是程序崩溃，也不是用户 ID 为 \code{-1}：

\begin{lstlisting}[language=C++,caption={特殊值会污染语义}]
int find_user_id(std::string_view name)
{
    if (name == "Ada") {
        return 1;
    }
    return -1;
}
\end{lstlisting}

\code{-1} 是约定，不是类型。调用者可能忘记检查，也可能把 \code{-1} 当成普通 ID 继续传。标准库词汇类型的价值，是把常见语义放进类型：可能没有值、若干类型之一、共享文本视图、连续区间视图。

\section{optional：可能没有值}

\code{std::optional<T>} 表达“有一个 \code{T}，或没有”：

\begin{lstlisting}[language=C++,caption={optional 表达查找结果}]
std::optional<User> find_user(std::string_view name)
{
    if (name == "Ada") {
        return User{.id = 1, .name = "Ada"};
    }
    return std::nullopt;
}
\end{lstlisting}

调用者必须面对没有值：

\begin{lstlisting}[language=C++,caption={检查 optional}]
const std::optional<User> user = find_user("Grace");
if (!user.has_value()) {
    std::cout << "not found\n";
    return;
}
std::cout << user->name << '\n';
\end{lstlisting}

\code{optional} 不适合表达复杂错误原因。如果调用者需要知道“文件不存在、格式错误、权限不足”，就需要错误类型。

\section{variant：若干类型之一}

命令行参数可以是不同形态：

\begin{lstlisting}[language=C++,caption={命令行值的 variant}]
struct Flag {
    std::string name;
};

struct KeyValue {
    std::string key;
    std::string value;
};

using Argument = std::variant<Flag, KeyValue>;
\end{lstlisting}

处理时用 \code{std::visit}：

\begin{lstlisting}[language=C++,caption={访问 variant 参数}]
void print_argument(const Argument& argument)
{
    std::visit([](const auto& item) {
        print_one(item);
    }, argument);
}
\end{lstlisting}

\code{variant} 适合类型集合固定的情况。它让编译器知道所有可能形态，也能在新增形态时提醒处理点。

\section{expected 风格：值或错误}

\cpp23 有 \code{std::expected}，但很多 \cpp20 工程会自己定义简单结果类型，或使用成熟库。核心思想是一样的：成功时有值，失败时有错误。

\begin{lstlisting}[language=C++,caption={expected 风格结果}]
template <typename T, typename E>
struct Result {
    std::optional<T> value;
    std::optional<E> error;
};
\end{lstlisting}

这个简化版本不能阻止同时有值和错误。更严谨可以用 \code{std::variant<T, E>}：

\begin{lstlisting}[language=C++,caption={variant 表达二选一}]
template <typename T, typename E>
class Result2 {
public:
    static Result2 success(T value)
    {
        return Result2{std::move(value)};
    }

    static Result2 failure(E error)
    {
        return Result2{std::move(error)};
    }

private:
    explicit Result2(std::variant<T, E> state)
        : state_{std::move(state)}
    {
    }

    std::variant<T, E> state_;
};
\end{lstlisting}

教材里经常用简单结构，是为了降低语法负担；工程里要根据风险选择更严谨的表示。

\section{pair 和 tuple 不要滥用}

\code{std::pair} 和 \code{std::tuple} 适合临时组合值，但不适合长期表达业务对象：

\begin{lstlisting}[language=C++,caption={tuple 让含义消失}]
std::tuple<std::string, int, bool> user;
\end{lstlisting}

读者不知道三个字段分别是什么。更好的写法是结构体：

\begin{lstlisting}[language=C++,caption={结构体保留字段含义}]
struct UserRow {
    std::string name;
    int age = 0;
    bool active = false;
};
\end{lstlisting}

如果返回两个迭代器、两个坐标、键值对，\code{pair} 合理；如果字段有业务名，就写类型。

\section{span 和 string\_view 也是词汇类型}

\code{std::span<const T>} 表达非拥有连续区间，\code{std::string_view} 表达非拥有文本视图。它们减少复制，但也把生命周期要求推给调用者。词汇类型不是越轻越好，而是要准确表达语义。

例如解析函数只观察文本：

\begin{lstlisting}[language=C++,caption={解析只需要观察文本}]
std::optional<int> parse_port(std::string_view text);
\end{lstlisting}

配置对象要保存文本：

\begin{lstlisting}[language=C++,caption={保存文本要拥有字符串}]
struct ServerConfig {
    std::string host;
};
\end{lstlisting}

不要把 \code{string_view} 存进长期配置对象，除非你清楚底层字符串的生命周期。

\section{完整案例：解析命令行参数}

一个参数可以是标志、键值对或位置参数：

\begin{lstlisting}[language=C++,caption={参数词汇类型}]
struct Positional {
    std::string value;
};

using ParsedArgument = std::variant<Flag, KeyValue, Positional>;

ParsedArgument parse_argument(std::string_view text)
{
    if (text.starts_with("--")) {
        const std::size_t equal = text.find('=');
        if (equal == std::string_view::npos) {
            return Flag{.name = std::string{text.substr(2)}};
        }
        return KeyValue{
            .key = std::string{text.substr(2, equal - 2)},
            .value = std::string{text.substr(equal + 1)},
        };
    }
    return Positional{.value = std::string{text}};
}
\end{lstlisting}

这个函数把不同形态显式化。后续处理参数时，编译器知道所有可能情况，而不是让一堆字符串约定散在代码里。

\begin{exercisebox}
把前面的 \code{ParseResult} 改成 \code{std::variant<Options, ParseError>} 风格。写一个辅助函数判断成功或失败。比较这种写法和两个 \code{optional} 字段的优缺点。
\end{exercisebox}

\section{optional 背后的状态模型}

\code{std::optional<T>} 可以用普通代码近似理解成“一个是否有值的标志，加上一块能放 \code{T} 的存储”。它不是把 \code{T} 初始化成某个特殊值。

\begin{lstlisting}[language=C++,caption={optional 的近似模型}]
template <typename T>
class OptionalLike {
public:
    [[nodiscard]] bool has_value() const noexcept
    {
        return this->has_value_;
    }

private:
    bool has_value_ = false;
    // 真实 optional 会用合适的未初始化存储保存 T
};
\end{lstlisting}

这解释了为什么 \code{optional<int>} 能区分“没有值”和“值为 \code{0}”。特殊值方案把状态塞进值域；\code{optional} 把状态单独拿出来。

\section{variant 背后的标签联合体}

\code{std::variant<Flag, KeyValue, Positional>} 可以理解成“一个标签，加上一块足够大的存储”。标签记录当前里面是哪一种类型：

\begin{lstlisting}[language=C++,caption={variant 的近似状态}]
enum class ArgumentKind {
    flag,
    key_value,
    positional,
};

struct ArgumentLike {
    ArgumentKind kind = ArgumentKind::flag;
    // 真实 variant 会管理构造、析构、对齐和异常安全
};
\end{lstlisting}

手写这种结构很容易漏析构、拷贝和异常安全，所以标准库提供 \code{variant}。但理解“标签加存储”很重要：处理 variant 时必须看标签，也就是用 \code{std::visit} 或 \code{std::get_if} 分支处理当前形态。

\section{visit 背后是按标签分发}

下面的 \code{visit}：

\begin{lstlisting}[language=C++,caption={visit 处理三种参数}]
std::visit([](const auto& item) {
    print_one(item);
}, argument);
\end{lstlisting}

可以近似理解成：

\begin{lstlisting}[language=C++,caption={按当前标签分发}]
switch (argument.kind()) {
case ArgumentKind::flag:
    print_one(argument.as_flag());
    break;
case ArgumentKind::key_value:
    print_one(argument.as_key_value());
    break;
case ArgumentKind::positional:
    print_one(argument.as_positional());
    break;
}
\end{lstlisting}

标准库帮你保证类型安全。新增一种参数形态时，访问器如果处理不了，编译器会在相应调用点报错。这比字符串里塞一个 \code{"kind"} 字段更早暴露遗漏。

\section{expected 风格要保证二选一}

两个 \code{optional} 字段的简化 \code{Result} 有一个设计漏洞：它可能同时有值和错误，也可能两者都没有。用 \code{variant<T, E>} 表达二选一更严谨，因为状态空间只允许成功或失败。

\begin{lstlisting}[language=C++,caption={二选一结果的最小接口}]
template <typename T, typename E>
class Result {
public:
    [[nodiscard]] bool has_value() const noexcept
    {
        return std::holds_alternative<T>(this->state_);
    }

    [[nodiscard]] const T& value() const
    {
        return std::get<T>(this->state_);
    }

    [[nodiscard]] const E& error() const
    {
        return std::get<E>(this->state_);
    }

private:
    std::variant<T, E> state_;
};
\end{lstlisting}

真实项目还要处理构造函数、移动、错误访问策略和易用性。本章重点是模型：如果业务语义是二选一，类型最好也只允许二选一。

\section{词汇类型选择表}

\begin{longtable}{p{0.25\linewidth}p{0.29\linewidth}p{0.36\linewidth}}
\toprule
语义 & 类型 & 不要误用成 \\
\midrule
可能没有值 & \code{std::optional<T>} & 特殊值 \code{-1}、空字符串 \\
固定几种形态之一 & \code{std::variant<A, B>} & 手写 \code{kind} 加裸字段 \\
值或错误 & expected 风格结果 & 打印错误后返回默认值 \\
非拥有连续区间 & \code{std::span<T>} & 指针加长度分离传参 \\
非拥有文本 & \code{std::string_view} & 长期保存的配置字段 \\
业务字段集合 & \code{struct} & 长期使用的 \code{tuple} \\
\bottomrule
\end{longtable}

词汇类型的价值是让读者在函数签名里看见语义。签名越诚实，调用者越不容易猜错。
